% ============================================================================
\section{Kết luận}
\label{sec:conclusion}
% ============================================================================

\subsection{Kết luận}

Đồ án khảo sát hai mô hình Transformer tiền huấn luyện cho đọc hiểu trích xuất tiếng Việt --- PhoBERT
(đơn ngữ, word-level) và XLM-R (đa ngữ, subword) --- được fine-tune trong cùng điều kiện trên
UIT-ViQuAD~2.0 và chấm một lần trên toàn bộ \NVal{} câu validation. Bốn câu hỏi nghiên cứu được trả lời
như sau.

\begin{enumerate}
  \item \textbf{RQ1 --- ranh giới từ.} Giả thuyết ghi trước (PhoBERT yếu hơn vì lỗi tách từ) \emph{sai}
        ở dạng tổng quát: PhoBERT hơn XLM-R \PairAnsDiff{} điểm EM trên câu có đáp án, và tỉ lệ lỗi biên
        khi đã trả lời của hai mô hình như nhau (\PhobertBoundaryRate\% và \XlmrBoundaryRate\%). Cơ chế
        ``không trả lời được nửa từ'' có thật, nhìn thấy được qua dấu hiệu F1 cao/EM thấp, nhưng chỉ tác
        động lên \SegMisPct\% câu hỏi và bị chặn trên ở \SegCeilingLoss{} điểm EM; phần lớn các câu đó là lỗi
        của bộ tách từ \code{pyvi}, không phải ranh giới từ ghép thật.
  \item \textbf{RQ2 --- độ dài.} Không có suy giảm theo độ dài ngữ cảnh ở cả hai họ mô hình; độ dài đáp án
        làm giảm EM (không giảm F1) như nhau ở cả hai.
  \item \textbf{RQ3 --- câu không có đáp án.} Không mô hình nào thật sự nhận ra câu không có đáp án. XLM-R từ
        chối đúng tần suất nhưng sai đối tượng (precision \XlmrAbsPrec\%); PhoBERT gần như không từ chối
        (\PhobertAbsRate\%), vì loss trên lớp ``không có đáp án'' của nó tăng khoảng 7 lần sau epoch 1 trong
        khi loss trên vị trí đáp án vẫn giảm --- lặp lại ở hai tốc độ học. Cả hai thường trả về đúng cụm
        gây nhiễu mà người gán nhãn đặt vào câu hỏi.
  \item \textbf{RQ4 --- baseline.} Transformer vượt TF-IDF ở mọi nhóm; TF-IDF thậm chí không vượt được hệ
        thống luôn từ chối về EM tổng.
\end{enumerate}

\paragraph{Điều rút ra.} Điểm tổng của PhoBERT và XLM-R không khác nhau có ý nghĩa thống kê, nhưng hai mô
hình có hai hồ sơ kỹ năng ngược nhau: một mô hình đọc giỏi mà không biết từ chối, một mô hình biết từ chối
nhưng từ chối nhầm chỗ. Chọn mô hình dựa trên điểm tổng sẽ bỏ qua điều quan trọng nhất cho triển khai. Với
câu hỏi ban đầu của môn học --- \emph{đặc điểm dữ liệu quyết định phương pháp} --- câu trả lời của đồ án là:
trên UIT-ViQuAD~2.0, đặc điểm quyết định không phải là tách từ mà là \emph{một phần ba câu hỏi không có đáp
án}; và việc một mô hình xử lý phần đó ra sao phải được đo riêng.

\paragraph{Về phương pháp.} Hai trong ba công cụ chẩn đoán có giá trị nhất của đồ án không cần chạy mô hình
nào: phép đo biên từ cho biết trước giới hạn trên của hiệu ứng mà giả thuyết RQ1 dự đoán, và phép kiểm định
bộ stress-test cho thấy bộ dữ liệu chẩn đoán do chính nhóm xây \emph{không đo được thứ nó tuyên bố đo}. Nếu
không kiểm định, bảng kết quả trên bộ đó sẽ xếp ``không làm gì'' đứng đầu mà không ai nhận ra vì sao.

\subsection{Giới hạn}
\label{sec:limitations}

Các giới hạn dưới đây quyết định phạm vi mà mọi kết luận ở trên có giá trị.

\begin{enumerate}
  \item \textbf{Một seed huấn luyện cho mỗi mô hình.} Kiểm định McNemar và bootstrap đo độ bất định do
        \emph{chọn câu hỏi}, không đo độ bất định do \emph{seed}. Với mô hình cỡ base trên SQuAD, chênh
        lệch giữa các seed có thể lên tới một vài điểm EM; một chênh lệch nhỏ hơn cỡ đó không nên được
        hiểu là khác biệt giữa hai \emph{kiến trúc}.
  \item \textbf{Không tìm kiếm siêu tham số.} Cả hai mô hình dùng cùng một cấu hình phổ biến; mỗi mô
        hình có thể đạt điểm cao hơn ở cấu hình riêng của nó. So sánh ở đây là ``cùng điều kiện'', không
        phải ``mỗi bên ở trạng thái tốt nhất''.
  \item \textbf{Bộ tách từ khác với lúc tiền huấn luyện.} PhoBERT được tiền huấn luyện trên văn bản tách
        bởi RDRSegmenter; đồ án dùng \code{pyvi} để không phụ thuộc Java. Lệch phân phối này có thể làm
        PhoBERT kém hơn khả năng thật của nó, và Mục~\ref{sec:alignment-res} cho thấy \code{pyvi} có những
        lỗi hệ thống (dính dấu chấm, nối nhầm âm tiết).
  \item \textbf{Validation đóng vai tập test.} Test chính thức ẩn nhãn nên không chấm được. Validation
        không được dùng cho bất kỳ quyết định nào trước khi chấm (Mục~\ref{sec:protocol}), nhưng kết quả
        không so sánh trực tiếp được với các con số công bố trên test của cuộc thi VLSP.
  \item \textbf{Giới hạn vị trí khác nhau là một phần của so sánh.} PhoBERT dùng 256 token, XLM-R 384
        token. Đây là ràng buộc kiến trúc chứ không phải lựa chọn, nhưng nó có nghĩa là ``PhoBERT vs XLM-R''
        không cô lập riêng biến tokenization.
  \item \textbf{Bộ stress-test không dùng được làm thước đo} (Mục~\ref{sec:audit}). Các nhóm hiện tượng
        E2--E5 vì vậy không có bằng chứng định lượng đáng tin cậy trong đồ án này; phân tích theo độ dài và
        theo loại lỗi trên validation chỉ thay thế được một phần.
  \item \textbf{Phân loại thủ công 33 câu lệch biên} do một người thực hiện, chưa đo độ đồng thuận giữa
        các người gán. Phân loại lỗi tự động (Mục~\ref{sec:taxonomy}) dựa trên quan hệ tập hợp token nên
        tái lập được, nhưng không phân biệt được lỗi ``hiểu sai câu hỏi'' với ``chọn nhầm thực thể cùng loại''.
  \item \textbf{Miền dữ liệu.} Toàn bộ dữ liệu là Wikipedia; kết luận không tự động mở rộng sang văn bản
        mạng xã hội, nơi lỗi tách từ được dự đoán còn nặng hơn.
\end{enumerate}

\subsection{Hướng phát triển}
\label{sec:future}

\begin{enumerate}
  \item \textbf{Hiệu chỉnh ngưỡng từ chối trên dev.} Ngưỡng $\tau$ trong công thức~(\ref{eq:null}) được cố định
        bằng 0. Chọn $\tau$ riêng cho từng mô hình trên tập dev (không phải validation) là cách rẻ nhất để
        thử khôi phục khả năng từ chối của PhoBERT; kết quả cần được chấm lại trên validation.
  \item \textbf{Tìm nguyên nhân gốc của việc PhoBERT bỏ lớp ``không có đáp án''}: nhiều seed, ghi lại chuẩn
        gradient theo loại nhãn, thử trọng số lớp, thử độ dài 256 cho XLM-R để tách biến ``độ dài cửa sổ''
        khỏi biến ``mô hình''.
  \item \textbf{Bộ tách từ.} Chạy lại PhoBERT với RDRSegmenter (VnCoreNLP) --- bộ tách từ dùng lúc tiền
        huấn luyện --- để đo phần hiệu năng bị mất do lệch bộ tách từ; và sửa luật dính dấu chấm của
        \code{pyvi}, nguyên nhân của \MisPeriod{} trong \SegMis{} câu lệch biên.
  \item \textbf{Xây lại bộ stress-test} theo quy trình có kiểm định tự động từ đầu: đáp án phải nằm trong
        ngữ cảnh (báo lỗi thay vì trả về 0), mỗi nhóm có đủ câu có đáp án làm đối chứng, và câu hỏi do người
        viết thay vì sinh từ khuôn mẫu; đo độ đồng thuận giữa người gán.
  \item \textbf{Nhiều seed} cho mọi mô hình để tách độ bất định do huấn luyện khỏi độ bất định do chọn câu
        hỏi.
\end{enumerate}
